% ===== 原文 pp.1–5（Lead：§1 引言 + §2 开头 + §2.1 复群上的 Higgs 丛） =====
\pagenumbering{arabic}

% [原文 p.1]
\section{引言}

Higgs 丛为数学与物理的众多分支提供了一个统一的结构。$G_{\C}$-Higgs 丛的 Dolbeault 模空间
$\mathcal{M}_{G_{\C}}$ 具有超凯勒结构；通过不同的复结构，它可以被理解为不同的模空间：
\begin{itemize}[leftmargin=1.4em,itemsep=1pt,topsep=2pt]
\item 经由非阿贝尔 Hodge 对应，该模空间作为实流形微分同胚于平坦联络的 de Rham 模空间
  [36,41,60,95,100]。
\item 经由 Riemann--Hilbert 对应，de Rham 模空间与曲面群表示的 Betti 模空间之间存在解析对应。
\end{itemize}
% [Tutor Note: 原文 arXiv 版的 2010 Mathematics Subject Classification 一栏写的是 "Primary ????; Secondary ????"，
%  作者未填写，本译本照原样保留并在译者说明中注明，不代为猜测分类号。]

% [原文 p.2]
通过这些对应，Higgs 丛同时表现为平坦联络与表示——这两类对象在当代数学中无处不在，并与理论物理
密切相关。例如：
\begin{itemize}[leftmargin=1.4em,itemsep=1pt,topsep=2pt]
\item 通过 Hitchin 纤维化，$\mathcal{M}_{G_{\C}}$ 给出了作为可积系统的超凯勒流形的例子 [61]，
  从而在物理中有重要应用，例如在 Gukov、Hitchin、Donagi、Pantev 与 Witten 的工作中
  （例如见 [37,38,40,54,55,63,103]）。
\item Hausel--Thaddeus [58] 把 Higgs 丛与镜像对称联系起来，并与 Donagi--Pantev [40] 一起把
  $\mathcal{M}_{G_{\C}}$ 呈现为 Calabi--Yau 流形镜像对称的基本例子，其几何与拓扑至今仍被研究。
\item Kapustin--Witten [69] 利用 Higgs 丛，通过镜像对称得到了几何 Langlands 对应的一个物理推导。
  不久之后，Ngô 在证明 [79] 中的基本引理时发现 Higgs 丛是关键要素。
\end{itemize}
% [Tutor Note: 原文此处 "Through these correspondances" 中 correspondances 为拼写错误（应为 correspondences）；
%  正文按原文排版，正确形式见本批注。]

在本讲义中，我们将聚焦于 Higgs 丛及其模空间，考察该理论的三个主要方面：
\begin{itemize}[leftmargin=1.4em,itemsep=1pt,topsep=2pt]
\item \textbf{Hitchin 纤维化。}我们将把 Hitchin 纤维化作为理解 Higgs 丛的工具来回顾，并描述
  Hitchin 在 [61] 中引入的阿贝尔化，以及如 [64] 中那样、对某些 Higgs 丛子空间出现的非阿贝尔化。
\item \textbf{膜的构造。}我们将构造并研究 Higgs 丛模空间中的膜族。特别地，膜可以通过在黎曼面与规范群上
  施加不同的结构而得到。我们还将考虑通过膜出现的几何结构，例如超多边形空间。
\item \textbf{对应关系。}最后，我们还将考虑 Higgs 丛之间以及膜之间的对应关系，它们与数学和物理的
  不同领域相关：例如镜像对称、通过群自同构产生的对应，以及 3-流形。
\end{itemize}

关于 Higgs 丛谱数据的更长介绍，读者可参阅 [93] 与 [89] 中的讲义。其他文献将在本讲义各处提及。

\para{文献说明。}
我们将突出所考虑的主要文献，以及所用方法被发展出来的确切位置。由于对相关领域的研究给出全面而
详尽的综述被证明非常困难，我们只在相关工作直接涉及讲义中提到的方法时才予以提及。读者应参阅
参考文献中的文献以获取进一步信息（例如见 [7,83,89] 中的参考文献）。

% [原文 p.3]
\section{Higgs 丛模空间的几何}

\begin{quote}
\kaishu 证明是创造性想象与批判性推理之间长期相互作用的最终产物。没有证明，纲领仍不完整；
但没有想象力的投入，它根本不会开始。
\end{quote}
\begin{flushright}
——迈克尔·阿蒂亚爵士（Sir Michael Atiyah）
\end{flushright}

我们将用本章回顾 Higgs 丛的基本概念，以及它们向主 $G_{\C}$-Higgs 丛、实 Higgs 丛、抛物 Higgs 丛
与野 Higgs 丛的若干推广（更多细节见 [36,41,60,61,80,95,96]）。

\subsection{复群上的 Higgs 丛}

除另有说明外，我们总设 $\Sigma$ 是亏格 $g\ge 2$ 的紧黎曼面，其典范丛为 $K=T^*\Sigma$。

\begin{defn}{2.1}
一个 \emph{Higgs 丛}是一对 $(E,\Phi)$，其中 $E$ 是 $\Sigma$ 上的全纯向量丛，而 \emph{Higgs 场}
$\Phi\in H^0(\Sigma,\End(E)\otimes K)$。
\end{defn}

在本文中，我们把定义 2.1 中的上述 Higgs 丛称为\emph{经典 Higgs 丛}。回忆：$\Sigma$ 上的全纯向量丛
$E$ 由其秩 $\rk(E)$ 与次数 $\deg(E)$ 作拓扑分类，由此可以定义其\emph{斜率}为
\[
  \mu(E) := \deg(E)/\rk(E).
\]
称向量丛 $E$ 是\emph{稳定}的（或\emph{半稳定}的），若对任意真、非零子丛 $F\subset E$ 都有
$\mu(F)<\mu(E)$（或 $\mu(F)\le\mu(E)$）。若 $E$ 是若干稳定丛的直和，且这些稳定丛的斜率均为 $\mu(E)$，
则称 $E$ 是\emph{多稳定}的。固定次数 $d$ 与秩 $n$ 的稳定丛的模空间是代数簇 $\mathcal{N}(n,d)$，
它可由 Mumford 的几何不变量理论构造。此外，当 $n$ 与 $d$ 互素时，该空间是维数为
$n^2(g-1)+1$ 的光滑射影代数簇。
% [Tutor Note: 注意两个模空间的记号不同：稳定向量丛的模空间是 $\mathcal{N}(n,d)$，而后面 Higgs 丛的
%  模空间是 $\mathcal{M}(n,d)$；$\mathcal{M}_{G_{\C}}$ 则表示主 $G_{\C}$-Higgs 丛的模空间。三者不要混。]

\begin{exmp}{2.2}
由于线丛都是稳定的，空间 $\mathcal{N}(1,d)$ 包含所有次数为 $d$ 的线丛，并且同构于雅可比簇
$\Jac^d(\Sigma)$——一个维数为 $g$ 的阿贝尔簇。
\end{exmp}

把向量丛的稳定性条件推广到 Higgs 丛，我们考虑如下定义：

\begin{defn}{2.3}
设 $F$ 是 $E$ 的向量子丛。若 $\Phi(F)\subset F\otimes K$，则称 $F$ 是 $E$ 的 \emph{$\Phi$-不变子丛}。
称 Higgs 丛 $(E,\Phi)$ 为
\begin{itemize}[leftmargin=1.6em,itemsep=2pt,topsep=2pt]
\item \emph{稳定}的（\emph{半稳定}的），若对每个真 $\Phi$-不变子丛 $F\subset E$ 都有
  $\mu(F)<\mu(E)$（等价地，$\mu(F)\le\mu(E)$）；
\item \emph{多稳定}的，若
  \[
    (E,\Phi) = (E_1,\Phi_1)\oplus(E_2,\Phi_2)\oplus\cdots\oplus(E_r,\Phi_r),
  \]
  其中 $(E_i,\Phi_i)$ 稳定，且对一切 $i$ 有 $\mu(E_i)=\mu(E)$。
\end{itemize}
\end{defn}

% [原文 p.4]
\begin{rem}{2.4}
$\Phi$ 限制在不变子丛上的特征多项式整除 $\Phi$ 的特征多项式。因此，若 Higgs 场的特征多项式不可约，
则相应的 Higgs 丛自动是稳定的。
\end{rem}
% [Tutor Note: 此处的论证要点是：若 $\Phi$ 在某个真不变子丛上的特征多项式是 $\Phi$ 特征多项式的真因子，
%  则特征多项式可约；反之不可约就排除了真不变子丛，故自动稳定（不需要额外的斜率条件）。]

\begin{exmp}{2.5}
取典范丛 $K$ 的一个平方根 $K^{1/2}$，以及 $K^2$ 的一个截面 $\omega$。考虑向量丛
$E=K^{\frac12}\oplus K^{-\frac12}$ 以及由
\[
  \Phi_\omega = \begin{pmatrix} 0 & \omega \\ 1 & 0 \end{pmatrix}
  \in H^0(\Sigma,\End(E)\otimes K)
\]
给出的 Higgs 场 $\Phi_\omega$，便得到一族经典 Higgs 丛 $(E,\Phi_\omega)$（如定义 2.1）。
\end{exmp}

考虑一个严格半稳定的 Higgs 丛 $(E,\Phi)$。由于 $E$ 不稳定，它有一个极小的 $\Phi$-不变子丛
$F\subset E$，使得 $\mu(F)=\mu(E)$。此时诱导对 $(F,\Phi)$ 稳定，而诱导商（对 $(E/F,\Phi')$）
半稳定。归纳地，可以得到一个子丛旗
\[
  F_0=0\subset F_1\subset\cdots\subset F_r=E,
\]
其中 $\mu(F_i/F_{i-1})=\mu(E)$（$1\le i\le r$），且诱导的 Higgs 丛 $(F_i/F_{i-1},\Phi_i)$ 稳定。
这就是 $E$ 的 \emph{Jordan--Hölder 滤链}；它虽不唯一，但诱导出一个在同构意义下唯一的 graded 对象
\[
  \mathrm{Gr}(E,\Phi) := \bigoplus_{i=1}^{r}(F_i/F_{i-1},\Phi_i).
\]
若 $\mathrm{Gr}(E,\Phi)\cong\mathrm{Gr}(E',\Phi')$，则称两个半稳定 Higgs 丛 $(E,\Phi)$ 与
$(E',\Phi')$ 是 \emph{$S$-等价}的。按照 [80, Theorem 5.10]，记 $\mathcal{M}(n,d)$ 为固定次数 $d$、
固定秩 $n$ 的半稳定 Higgs 丛的 $S$-等价类的模空间。模空间 $\mathcal{M}(n,d)$ 是拟射影概形，
并有一个开子概形 $\mathcal{M}'(n,d)$，它是稳定对的模概形。于是每一点都由一个稳定或多次稳定的
Higgs 丛代表。当 $d$ 与 $n$ 互素时，模空间 $\mathcal{M}(n,d)$ 是复维数为 $2n^2(g-1)+2$ 的光滑
非紧簇，且 $\mathcal{N}(n,d)$ 在稳定轨迹上的余切空间作为 Zariski 开子集包含在 $\mathcal{M}(n,d)$ 中。
% [Tutor Note: 原文 "the induced quotient the quotient $(E/F,\Phi')$" 中 the quotient 重复，属原文笔误；
%  正确应读作 "the induced quotient $(E/F,\Phi')$"。另：$\dim\mathcal{M}(n,d)=2n^2(g-1)+2$
%  恰好是 $\dim\mathcal{N}(n,d)=n^2(g-1)+1$ 的两倍，这来自 $\mathcal{M}(n,d)$ 在稳定轨迹上与
%  $T^*\mathcal{N}(n,d)$ 的 Zariski 开子集重合——可用于验算。]

一个最重要的刻画来自 Hitchin [61] 与 Simpson [95] 的工作，并可以推广到更一般的情形：

\begin{thm}{2.6}
若 Higgs 丛 $(E,\Phi)$ 稳定且 $\deg E=0$，则 $E$ 上存在唯一的酉联络 $A$（与全纯结构相容），使得
\begin{equation}
  F_A + [\Phi,\Phi^*] = 0 \in \Omega^{1,1}(\Sigma,\End E), \tag{2.7}
\end{equation}
其中 $F_A$ 是联络的曲率。
\end{thm}

式 (2.7) 与全纯性条件
\begin{equation}
  d''_A\Phi = 0 \tag{2.8}
\end{equation}

% [原文 p.5]
一起称为 \emph{Hitchin 方程}，其中 $d''_A\Phi$ 是 Higgs 场 $\Phi$ 的协变导数的反全纯部分。此外，
Hitchin 还证明了 Higgs 丛的模空间是超凯勒流形，其自然辛形式 $\omega$ 定义在 Higgs 丛的无穷小形变
$(\dot A,\dot\Phi)$ 上：
\[
  \omega\big((\dot A_1,\dot\Phi_1),(\dot A_2,\dot\Phi_2)\big)
  = \int_\Sigma \tr(\dot A_1\dot\Phi_2-\dot A_2\dot\Phi_1),
\]
其中 $\dot A\in\Omega^{0,1}(\End_0 E)$，$\dot\Phi\in\Omega^{1,0}(\End_0 E)$（详见 [60,61]）。
为简便起见，我们固定 $n$ 与 $d$，并记 $\mathcal{M}$ 为 $\mathcal{M}(n,d)$。
% [Tutor Note: 辛形式公式里的 $\dot A$ 落在 $\Omega^{0,1}$、$\dot\Phi$ 落在 $\Omega^{1,0}$，两者都取
%  $\End_0 E$（迹零自同态）系数；积分 $\int_\Sigma$ 需要 $\dot A_1\dot\Phi_2$ 这类 $(1,1)$-形式才可积，
%  这正是取 $\Omega^{0,1}$ 与 $\Omega^{1,0}$ 的原因。]

Higgs 丛的概念可以推广到主 $G_{\C}$-丛，其中 $G_{\C}$ 是复半单李群。更多细节读者可参阅 [61]。

\begin{defn}{2.9}
一个 \emph{$G_{\C}$-Higgs 丛}是一对 $(P,\Phi)$，其中 $P$ 是 $\Sigma$ 上的主 $G_{\C}$-丛，
而 Higgs 场 $\Phi$ 是向量丛 $\mathrm{ad}P\otimes_{\C} K$ 的全纯截面，这里 $\mathrm{ad}P$ 是伴随表示
所对应的向量丛。
\end{defn}

当 $G_{\C}\subset\Gl(n,\C)$ 时，$G_{\C}$-Higgs 丛给出经典意义下的 Higgs 丛，并带有反映 $G_{\C}$
定义的额外结构。特别地，经典 Higgs 丛正是 $\Gl(n,\C)$-Higgs 丛。
% [Tutor Note: 这里 "$\mathrm{ad}P$" 是伴随丛（adjoint bundle）$P\times_{\mathrm{Ad}}\mathfrak{g}_{\C}$，
%  不是某个算子的核；原文用 adP 而非 $\operatorname{ad}P$ 只是字体差异。]

\begin{exmp}{2.10}
例 2.5 中的 Higgs 丛其 Higgs 场无迹，且行列式 $\Lambda^2E$ 平凡。因此，对每个二次微分 $\omega$
都有一个 $\Sl(2,\C)$-Higgs 丛 $(E,\Phi_\omega)$。
\end{exmp}

设 $G_{\C}$ 为复半单李群，记 $\mathcal{M}_{G_{\C}}$ 为多次稳定 $G_{\C}$-Higgs 丛的 $S$-等价类的模空间。
按照 [84]，可以对主 $G_{\C}$-丛定义稳定性；当考虑 $A,B,C,D$ 型群时，这可以用经典 Higgs 丛的稳定性
来表达（系统的研究见 [7, Section 1.1]）。此外，多次稳定的概念也可以搬到主 $G_{\C}$-丛上，从而可以
在紧黎曼面 $\Sigma$ 上构造固定拓扑类型的多次稳定主 $G_{\C}$-丛的同构类的模空间。由于在本讲义中
我们将处理自动稳定的 Higgs 对（如注 2.4），我们不再花时间回顾主 Higgs 丛稳定性的主要研究。
关于相应构造的细节，读者可参阅例如 [67, Section 3] 与 [7, Section 1]。

\begin{exmp}{2.11}
一个 $\SO(2n,\C)$-Higgs 丛是一对 $(E,\Phi)$，其中 $E$ 是秩为 $2n$ 的全纯向量丛，带有一个非退化
对称双线性型 $(\cdot,\cdot)$，而 $\Phi\in H^0(\Sigma,\End_0(E)\otimes K)$ 是满足
\begin{equation}
  (\Phi v,w) = -(v,\Phi w) \tag{2.12}
\end{equation}
的 Higgs 场。等价地，一个 $\SO(2n+1,\C)$-Higgs 丛是一对 $(E,\Phi)$，其中 $E$ 是秩为 $2n+1$ 的
全纯向量丛，带有非退化对称双线性型 $(v,w)$，而 $\Phi$ 是 $H^0(\Sigma,\End_0(E)\otimes K)$ 中满足
式 (2.12) 的 Higgs 场。
\end{exmp}
% [Tutor Note: 原文两处写法不统一：$(\Phi v,w)=-(v,\Phi w)$ 与 $(v,w)$ 都用了同一记号表示该双线性型；
%  注意 $\End_0$ 表示迹零自同态，对 $\SO$ 型这才是对的（$\Phi$ 取值在正交李代数中）。]
